\section{Bounded physical leaves and arbitrary chain lengths}
\label{sec:ldp_unequal_mera}

The isometries in the network need not be equal within a layer.
Allowing unequal positive block lengths removes the divisibility restriction,
while a uniform bound on the physical leaves keeps every local dimension
independent of the chain length and the requested accuracy.

\begin{definition}[Isometry trees with bounded leaves]
    \label{def:ldp_bounded_isometry_tree}
    \lean{MPSPreparation.IsometryTree, MPSPreparation.IsometryTree.matrix,
        MPSPreparation.IsometryTree.matrix_fork_apply,
        MPSPreparation.IsometryTree.cast, MPSPreparation.IsometryTree.matrix_cast}
    \leanok
    Fix a physical dimension $d$, a virtual dimension $\chi$ and a leaf
    bound $R$. A depth-zero tree on $n$ sites, with $1\le n\le R$, is
    an isometry $V\colon\C^\chi\to(\C^d)^{\otimes n}$.
    A tree of coarse depth $h+1$ consists of two trees $T_1,T_2$ of
    coarse depth $h$, on $n_1,n_2$ sites, and an isometry
    $W\colon\C^\chi\to\C^\chi\otimes\C^\chi$.
    Its contraction is
    \begin{align}
        T & =(T_1\otimes T_2)W
        \colon\C^\chi\longrightarrow(\C^d)^{\otimes(n_1+n_2)}.
        \label{eq:ldp_unequal_tree_contraction}
    \end{align}
    The physical sites keep their order under concatenation.
    A tree of coarse depth $h$ has $2^h$ leaves and $h+1$ isometry
    layers, counting the finest layer.
\end{definition}

\begin{theorem}[Isometric contraction of a tree]
    \label{thm:ldp_bounded_tree_isometry}
    \lean{MPSPreparation.IsometryTree.isIsometry_matrix}
    \leanok
    \uses{def:ldp_bounded_isometry_tree}
    Every tree contraction $T$ satisfies $T^\dagger T=\Id$.
\end{theorem}
\begin{proof}
    \leanok
    Induct on the tree. At a binary vertex,
    $T^\dagger T=W^\dagger(T_1^\dagger T_1\otimes T_2^\dagger T_2)W
        =W^\dagger W=\Id$.
\end{proof}

\begin{theorem}[Polar trees for unequal block lengths]
    \label{thm:ldp_unequal_polar_tree}
    \lean{MPSPreparation.exists_isometryTree_matrix_eq_cfgPolarIso}
    \leanok
    \uses{def:ldp_bounded_isometry_tree, def:ldpolar_matrix, def:injective}
    Let $s\ge1$, and assume that $A$ blocked over every $m\ge s$ sites
    is injective. If $2^h s\le n\le2^h c$, its polar isometry $V_n$
    is the contraction of a tree of coarse depth $h$ whose physical
    leaves contain at most $c$ sites.
\end{theorem}
\begin{proof}
    \leanok
    \uses{lem:ldp_polar_merge}
    Split $n$ into $n_1=\lfloor n/2\rfloor$ and $n_2=n-n_1$.
    At a successor depth both satisfy
    $2^h s\le n_i\le2^h c$; no rounding error is added to the upper
    bound, since its endpoint is even. The polar-merge identity gives
    $V_{n_1+n_2}=(V_{n_1}\otimes V_{n_2})W_{n_1,n_2}$.
    Injectivity makes $W_{n_1,n_2}$ an isometry. Induction gives the
    two descendants, and at depth zero the polar isometry is itself
    the bounded leaf.
\end{proof}

\begin{definition}[Finite-range networks on unequal blocks]
    \label{def:ldp_unequal_mera}
    \lean{MPSPreparation.UnequalMERA, MPSPreparation.UnequalMERA.topPair,
        MPSPreparation.UnequalMERA.treeMatrix,
        MPSPreparation.UnequalMERA.treeMatrix_apply,
        MPSPreparation.UnequalMERA.state, MPSPreparation.UnequalMERA.state_apply,
        MPSPreparation.UnequalMERA.ofTrees}
    \leanok
    \uses{def:ldp_bounded_isometry_tree}
    Let $D,M\ge1$, let $\ell_0+\cdots+\ell_{M-1}=N$, and let $T_k$
    be trees of common coarse depth $h$, virtual dimension $D^2$ and
    physical leaf bound $R$, on their respective blocks of $\ell_k$
    sites. A unitary $u$ on $\C^D\otimes\C^D$ prepares a pair
    $\ket\omega=u\ket{0,0}$. The network state is
    \begin{align}
        \ket\psi
         & =\left(\bigotimes_{k=0}^{M-1}T_k\right)
        \bigotimes_{k=0}^{M-1}\ket\omega_{R_kL_{k+1}},
        \label{eq:ldp_unequal_mera_state}
    \end{align}
    with cyclic indexing of the pair legs. All lower disentanglers
    are identities. The network has $h+2$ layers, including the finest
    isometries and the top unitary layer. Internal tree legs have
    dimension $D^2$; every physical leaf has dimension at most $d^R$.
\end{definition}

\begin{theorem}[Normalization and arbitrary top pairs]
    \label{thm:ldp_unequal_mera_norm}
    \lean{MPSPreparation.UnequalMERA.isIsometry_treeMatrix,
        MPSPreparation.UnequalMERA.sum_star_topPair_mul,
        MPSPreparation.UnequalMERA.norm_state,
        MPSPreparation.UnequalMERA.exists_disentangler_pair,
        MPSPreparation.UnequalMERA.exists_of_trees}
    \leanok
    \uses{def:ldp_unequal_mera}
    The state in (\ref{eq:ldp_unequal_mera_state}) has norm one.
    Conversely, for any such forest and any unit pair $\ket\omega$,
    a top unitary realizes that pair and hence the forest contraction.
\end{theorem}
\begin{proof}
    \leanok
    \uses{thm:ldp_bounded_tree_isometry, thm:ldp_chain_block_iso_inner,
        thm:ldp_isometric_extension_overlap, lem:ldp_unitary_implementing_isometry}
    Complete the unit column $\ket\omega$ to a unitary taking
    $\ket{0,0}$ to it. Tensor products of the normalized pairs have
    norm one, and the isometric forest preserves their norm.
\end{proof}

\begin{theorem}[Polar forests and the block-isometry state]
    \label{thm:ldp_unequal_mera_polar_state}
    \lean{MPSPreparation.UnequalMERA.state_eq_blockIsometryState,
        MPSPreparation.UnequalMERA.exists_of_trees_eq_blockIsometryState}
    \leanok
    \uses{def:ldp_unequal_mera, def:ldp_block_isometry_state}
    If $T_k=V_{\ell_k}$ and the top unitary prepares $\ket\omega$,
    the network state equals
    $(\bigotimes_k V_{\ell_k})\bigotimes_k\ket\omega_{R_kL_{k+1}}$.
    In particular, every normalized top pair and every forest of
    these polar trees give a network representing this block-isometry state.
\end{theorem}
\begin{proof}
    \leanok
    \uses{def:ldp_unequal_mera, def:ldp_block_isometry_state, thm:ldp_unequal_mera_norm}
    Substitute the tree contractions into
    (\ref{eq:ldp_unequal_mera_state}); the regrouping of physical
    configurations is precisely the prescribed block partition.
\end{proof}

\begin{theorem}[Bounded single leaves and product states]
    \label{thm:ldp_unequal_mera_caps}
    \lean{MPSPreparation.UnequalMERA.exists_single_leaf,
        MPSPreparation.UnequalMERA.nonempty_bond_one}
    \leanok
    \uses{def:ldp_unequal_mera}
    Every unit state on $1\le N\le R$ sites is represented exactly
    by a single leaf with virtual dimension one and coarse depth zero.
    If $d,R,N\ge1$, a product-state network of coarse depth zero
    also exists, with one physical site at each leaf and virtual
    dimension one.
\end{theorem}
\begin{proof}
    \leanok
    \uses{thm:ldp_unequal_mera_norm}
    For the first assertion use the unit column $1\mapsto\ket\psi$
    and the scalar top unitary. For the second use the column
    $1\mapsto\ket0$ separately at every site.
\end{proof}

\begin{theorem}[Uniform finite-range networks at every positive length]
    \label{thm:ldp_every_length_mera}
    \lean{MPSPreparation.exists_unequalMERA_every_length}
    \leanok
    \uses{def:ldp_unequal_mera, def:normal}
    Let $A$ be a normal tensor of physical dimension $d$ and bond
    dimension $D$. There are constants $R\ge1$ and $C\ge2$, chosen
    before $\varepsilon$ and $N$, such that for every $\varepsilon>0$
    and every positive length $N$ with nonzero periodic state, a
    network with some $1\le D'\le D$, physical leaf bound $R$ and
    coarse depth $h$ satisfies
    \begin{align}
        h+2
         & \le C\left[1+\log\!\left(1+\max\{0,\log(N/\varepsilon)\}\right)\right],
        \label{eq:ldp_every_length_mera_layers}                                    \\
        1-\left|\braket{\psi}{\phi_N(A)}\right|
         & \le\varepsilon.
        \label{eq:ldp_every_length_mera_error}
    \end{align}
    Its internal tree dimensions are at most $D^2$, and its physical
    leaf dimensions are at most $d^R$. Thus the layer count is
    $O(\log\log(N/\varepsilon))$ in the large-size, small-error regime,
    with no divisibility restriction on $N$.
\end{theorem}
\begin{proof}
    \leanok
    \uses{thm:ldp_unequal_polar_tree, thm:ldp_unequal_mera_polar_state,
        thm:ldp_unequal_mera_caps, thm:ldp_block_approximation_error,
        thm:ldp_short_chain_preparation, lem:ldp_normalized_state_rescaling,
        thm:ldp_depth_upper_bound_divisible, thm:prep_pair_norm}
    For $\varepsilon>1$, a product state suffices. Hence assume
    $0<\varepsilon\le1$. Pass to a normal left-canonical representative $B$, changing
    normalized states only by a phase. Choose $s\ge1$ from its
    eventual injectivity threshold and set $R=4s$. The unequal-block
    second-order estimate at $\gamma=1/2$ gives constants $K,r>0$
    with error at most $KM e^{-rq}$ for blocks of length at least $q$.
    Choose $q=\lceil a\log(N/\varepsilon)+b\rceil$, with $a=1/r$
    and fixed $b$ large enough that $q\ge s$ and $KN e^{-rq}\le\varepsilon$.

    If $q\le N$, take $M=\lfloor N/q\rfloor$ blocks, the first
    $M-1$ of length $q$ and the last of length $q+N\bmod q$.
    Choose $h$ so that $2^h s\le q<2^h(2s)$. Every block then lies
    between $2^h s$ and $2^h(4s)$, so
    Theorem~\ref{thm:ldp_unequal_polar_tree} supplies its bounded-leaf
    tree. Use the fixed-point pair at the top and apply the
    unequal-block error estimate.

    If $s\le N<q$, use one polar tree over $N$, with $h$ chosen
    from $N/s$. The normalized trace pair at its root gives the
    normalized periodic state exactly. This is not one leaf of
    accuracy-dependent width: its terminal widths remain at most $4s$.
    If $N<s$, the bounded single leaf of
    Theorem~\ref{thm:ldp_unequal_mera_caps} gives the state exactly,
    with $D'=1$.

    In the nontrivial branches $2^h\le q\le a\log(N/\varepsilon)+b+1$.
    Taking logarithms and adding the finest and top layers gives
    (\ref{eq:ldp_every_length_mera_layers}). The regularized expression
    also covers $N=1$ and $\varepsilon\ge1$.
\end{proof}

\begin{theorem}[Normal states at every sufficiently large length]
    \label{thm:ldp_eventual_every_length_mera}
    \lean{MPSPreparation.exists_unequalMERA_of_normal}
    \leanok
    \uses{def:ldp_unequal_mera, def:normal}
    For a normal tensor with $D\ge1$, the same fixed dimension and
    layer bounds hold for all $N\ge N_0$, for some fixed $N_0$,
    without a separate nonvanishing hypothesis.
\end{theorem}
\begin{proof}
    \leanok
    \uses{thm:ldp_every_length_mera, lem:ldp_periodic_state_nonzero}
    Normality implies eventual nonvanishing of the periodic state.
    Apply Theorem~\ref{thm:ldp_every_length_mera} beyond that threshold.
\end{proof}
